% Generated draft: requires mathematical and source review.
\documentclass[11pt]{article}
\usepackage[margin=25mm]{geometry}
\usepackage{fontspec}
\setmainfont{Latin Modern Roman}
\usepackage{amsmath,amssymb,mathtools,unicode-math}
\setmathfont{Latin Modern Math}
\usepackage{enumitem,needspace,xurl,hyperref}
\setlength{\parindent}{0pt}
\setlength{\parskip}{5pt}
\setlength{\emergencystretch}{3em}
\newcommand{\N}{\mathbb{N}}
\newcommand{\Z}{\mathbb{Z}}
\newcommand{\Q}{\mathbb{Q}}
\newcommand{\R}{\mathbb{R}}
\newcommand{\C}{\mathbb{C}}
\newcommand{\F}{\mathbb{F}}
\newcommand{\A}{\mathbb{A}}
\newcommand{\PP}{\mathbb{P}}
\newcommand{\E}{\mathbb{E}}
\newcommand{\eps}{\varepsilon}
\newcommand{\dd}{\,\mathrm{d}}
\begin{document}
\section*{UnsolvedMath ID: 10400051}
\textbf{Source problem code:} catalogue entry
\subsection{Definitions and mathematical statement}
For an oriented knot \(K\), the Kontsevich invariant \(Z(K)\) is a universal finite-type invariant valued in a degree-completed diagram space. Write the source's target as
\[
 \widehat{\mathcal A}(S^1)=\prod_{d\geq0}\mathcal A^{(d)}(S^1),\qquad
 Z(K)=\sum_{d\geq0}Z_d(K),
\]
where the graded pieces are diagram spaces on an oriented circle with their specified relations. In a Jacobi-diagram presentation these include antisymmetry, IHX and STU; the coefficient and framing conventions must match the source's \(Z\).

Problem 3.1 asks: for every oriented knot \(K\), calculate \(Z_d(K)\) for every degree \(d\), equivalently determine its full completed invariant \(Z(K)\). The task is an all-degrees computation, not merely a selected low-degree truncation. Since the full object is an infinite series, an exact all-degrees answer needs a formula or other specified effective degree-by-degree description, rather than an unspecified finite list of coefficients.

The excerpt does not fix a normalization, coefficient field, associator/construction convention, treatment of framing, or diagram basis/presentation for the output. For an unframed knot interpretation, the appropriate framing correction or quotient must be stated. It also does not prescribe a closed-form target, algorithmic input model, or complexity bound. These choices cannot be supplied by declaring the formal existence of the Kontsevich invariant to be the requested calculation.
\subsection{Short English statement}
Determine every homogeneous degree of the Kontsevich invariant of each oriented knot. The source asks for an all-degrees calculation, while leaving output format, normalization and computational criteria unspecified.
\subsection{Sources}
\begin{itemize}
\item Local UnsolvedMath registry, permanent ID 10400051, source code catalogue entry; no proposer or proposal year supplied; update 2026-08-06.
\item Ohtsuki, \emph{Problems on invariants of knots and 3-manifolds} (2002), Problem 3.1, PDF page 52: \url{https://msp.org/gtm/2002/04/gtm-2002-04-024s.pdf}. Registry access date 2026-07-29; PDF font-structure extraction. Provenance distinguishes 231 numbered bold headings. Current primary-literature status and redistribution terms require review.
\item Kontsevich is the supplied invariant attribution. The research summary gives no exact additional reference or resolution, describing the all-degrees task as a fundamental computational problem.
\end{itemize}
\end{document}
